% =============================================================
% Оборот титула (УДК, ББК, аннотация)
% =============================================================
% НЕ используем \thispagestyle{empty} — номер страницы должен быть (это страница 2)

\noindent УДК 621.891:532.516.5\\
ББК 34.41\\
З-14

\vspace{1cm}

\noindent\textbf{Авторы:}\\[0.3em]
\begin{center}
\textit{ассистент кафедры <<Технологические машины и~оборудование нефтегазового комплекса>> федерального государственного автономного образовательного учреждения высшего образования <<Сибирский федеральный университет>> А.~В.~Загуляев}\\[0.3em]
\textit{старший преподаватель кафедры <<Технологические машины и~оборудование нефтегазового комплекса>> федерального государственного автономного образовательного учреждения высшего образования <<Сибирский федеральный университет>> К.~А.~Башмур}
\end{center}

\vspace{0.5cm}

\noindent\textbf{Рецензенты:}\\[0.3em]
\begin{center}
доктор технических наук, профессор \textit{Е.~Д.~Агафонов}\\
доктор технических наук, профессор \textit{И.~В.~Ковалев}
\end{center}

\vspace{1cm}

Учебное пособие посвящено моделированию и~расчёту смазочных слоёв в~узлах трения с~дискретным микрорельефом. Рассмотрены конструкция и~принцип работы подшипников скольжения, постановки гидродинамических задач смазочного слоя на~основе уравнения Рейнольдса и~численные методы их решения с~учётом кавитации, термовязкостных эффектов и~шероховатости. Изложены методики расчёта радиальных и~упорных подшипников скольжения с~микрорельефом поверхности. Каждый раздел содержит контрольные вопросы для~самопроверки. Приведены примеры расчётов подшипниковых узлов и~задания для~курсового проектирования. Предназначено для~студентов направлений подготовки 15.03.02 и~15.04.02 <<Технологические машины и~оборудование>>, изучающих дисциплины <<Трение, износ и~смазка в~машинах>>, <<Технология повышения износостойкости объектов нефтегазового комплекса>>, <<Спецглавы механики жидкости и~газа>> и~другие. Соответствует требованиям образовательных стандартов.

\vfill

\noindent ISBN ХХХ-Х-ХХХХХ-ХХХ-Х

\vspace{0.3cm}

\begin{flushright}
\copyright~Загуляев А.~В., Башмур К.~А., \the\year\\
\copyright~Издательство <<Инфра-Инженерия>>, \the\year
\end{flushright}

\newpage
